\chapter{Consideraciones éticas y legales}

Las herramientas de detección de intrusiones operan sobre tráfico de red que puede contener información sensible, y un sistema que decide entre permitir y bloquear conlleva consecuencias operativas directas. Aunque esta tesis se limita a un prototipo experimental de inferencia offline, las decisiones de diseño que la hacen reproducible y segura tienen también una dimensión ética y legal que conviene hacer explícita.

\section{Uso dual y alcance estrictamente defensivo}
La investigación en ciberseguridad tiene una naturaleza intrínsecamente dual: las mismas técnicas que permiten detectar tráfico malicioso pueden, en principio, informar su evasión. El trabajo aquí descrito es defensivo por construcción y, además, deliberadamente inerte respecto a cualquier red real. Las salidas \texttt{PERMIT} y \texttt{BLOCK} son clasificaciones experimentales sobre registros de flujo estáticos: el sistema no realiza bloqueo en vivo, descarte de paquetes ni manipulación de \texttt{iptables}, y no se integra con cortafuegos, controladores de redes definidas por software (SDN) ni sistemas de prevención de intrusiones (IPS), tal como se delimita en los no objetivos del Capítulo~2. Esta ausencia de actuación en línea elimina el riesgo de que el artefacto, por sí mismo, provoque un bloqueo indebido de tráfico legítimo en una red en producción. El sistema es un instrumento de medición experimental, no un mecanismo de defensa desplegable.

\section{Privacidad y tratamiento de los datos}
El sistema opera exclusivamente sobre características estadísticas a nivel de flujo y no sobre el contenido (\emph{payload}) de los paquetes, evitando la inspección profunda de paquetes y buena parte de sus implicaciones legales y de privacidad. El contrato canónico de características refuerza esta posición al excluir a nivel de esquema los identificadores potencialmente sensibles o reidentificadores —direcciones IP de origen y destino, marcas temporales absolutas, \texttt{Flow ID} y números de puerto—, que se descartan antes de cualquier entrenamiento. El benchmark CICIDS2017 es un conjunto de datos público publicado por el Canadian Institute for Cybersecurity \parencite{Sharafaldin2018CICIDS2017,CICIDS2017Web}. El tráfico empleado en la Fase~2 fue generado por el propio operador en un laboratorio doméstico cerrado y no adversarial: se trata de tráfico propio, no contiene datos personales de terceros y no fue capturado de redes ajenas, por lo que su uso no plantea problemas de consentimiento ni de exposición de información personal.

\section{Riesgo operativo y asimetría de errores}
Las consecuencias de los dos tipos de error son profundamente asimétricas. Un falso negativo —un ataque clasificado como benigno y permitido— puede habilitar exfiltración de datos o interrupción del servicio; un falso positivo —tráfico benigno bloqueado— interrumpe comunicaciones legítimas y consume tiempo de análisis. El propio diseño de la función de recompensa codifica esta asimetría penalizando los falsos negativos con mayor severidad. Esta tesis reporta estos riesgos de forma honesta en lugar de ocultarlos tras la exactitud agregada: sobre la partición aleatoria fija \emph{in-distribution}, el modelo oficial MAIN presenta una tasa de falsos positivos de $0.00658$ y una tasa de falsos negativos de $0.00465$; sin embargo, bajo la partición por días, sustancialmente más exigente (Check~C), el recall de ataque desciende hasta $0.52954$. Esta caída del rendimiento ante un cambio de condiciones es precisamente la razón por la que el sistema se presenta como un prototipo de investigación y no como una herramienta lista para el despliegue: un uso operativo responsable exigiría validación sobre tráfico de producción etiquetado, múltiples entornos de captura y un análisis de seguridad independiente antes de delegar en el modelo cualquier decisión con efectos reales \parencite{Apruzzese2022CrossEvaluationNIDS,SommerPaxson2010ClosedWorld}.

\section{Reproducibilidad como práctica responsable}
Por último, la disciplina de reproducibilidad descrita en la metodología es también una práctica ética: la investigación en seguridad con aprendizaje automático adolece con frecuencia de código ausente y evaluaciones irreproducibles, lo que dificulta verificar o refutar las afirmaciones publicadas \parencite{Olszewski2023ReproSecurityML,Arp2020DosDontsMLSecurity}. La persistencia de artefactos por ejecución, la partición de prueba fija con semilla~42 y su manifiesto de referencia permiten que cualquier tercero audite exactamente bajo qué condiciones se obtuvo cada cifra reportada en esta tesis. Hacer las limitaciones verificables, en lugar de presentarlas como una declaración de buenas intenciones, es la forma de honestidad metodológica que este trabajo persigue.
